\documentclass[conference]{IEEEtran}

\usepackage{amsmath}
\usepackage{booktabs}
\usepackage{array}
\usepackage{url}

\title{Tabular Reinforcement Learning for ``The Floor Is Slime!''}

\author{
\IEEEauthorblockN{Alejandro Leon Martin Lama}
\IEEEauthorblockA{
Artificial Intelligence 2025/26\\
Universidad de Sevilla\\
Seville, Spain\\
alemarlam@alum.us.es}
}

\begin{document}
\maketitle

\begin{abstract}
This report presents a reinforcement learning solution for the long-jumps variant
of the two-player race game ``The floor is slime!''. The problem was modeled as a
finite stochastic decision process in which an agent must decide whether to stay,
jump one tile, or jump two tiles while accounting for changing terrain, stochastic
landing outcomes, and an opponent that also tries to reach the goal. The learning
algorithm is tabular Q-learning, implemented from scratch without reinforcement
learning libraries. The project includes a rule-based game engine, four dummy
opponent strategies, a training and evaluation pipeline, policy persistence, and
reproducible CSV/JSON experiment outputs. Particular attention was given to design
choices that affected learning quality: the state representation, reward shaping,
turn and round ownership, the use of a maximum round limit, and the choice of
training opponents. Experiments show that a compact local state representation
learns substantially better on the default 10-tile game than a full slime-mask
state, reaching a mean win rate of 0.578 after 20,000 episodes compared with a
0.411 random baseline. The full state is more theoretically Markovian, but it
suffers from severe state-space growth. On a 20-tile track, the full-state table
visited more than five million states. The main conclusion is that tabular
Q-learning is suitable for the assignment-sized game when supported by a compact
state abstraction and carefully balanced rewards, but its scalability is limited.
\end{abstract}

\begin{IEEEkeywords}
reinforcement learning, Q-learning, Markov decision process, tabular learning,
stochastic games, reward shaping
\end{IEEEkeywords}

\section{Introduction}

Reinforcement learning (RL) studies how an agent can learn to act by interacting
with an environment and receiving rewards or penalties \cite{sanchoQLearning}.
Unlike supervised learning, the agent is not given the exact correct answer after
an error. Instead, it receives an evaluation of its behavior and must discover a
policy that maximizes long-term return. This is appropriate for games because an
action that looks good immediately may be bad later, and an action that is risky
may still be necessary to win.

The selected assignment game is ``The floor is slime!''. It is a two-player race on
a line of tiles. Each player starts at the first tile, and the first player to
reach the final tile wins. Interior tiles may be dry or slimy, and after every
round empty tiles may flip between both terrain types. The selected variant is the
long-jumps version. In this version a player may stay, jump one tile, or attempt a
two-tile long jump. The long jump is strategically important because it can win
races quickly, but it has a lower landing probability, especially when the target
tile is slimy.

The objective of this project is to implement a complete RL pipeline for this game
from first principles, following the requirements of the course assignment
\cite{assignment2025}. The assignment requires an original algorithmic
implementation, a simulated dummy player for training, a way for a user to play
against computer-controlled behavior, reproducible experiments, and a scientific
report. The implementation therefore focuses on clarity and explainability rather
than on neural-network scale. A tabular Q-learning method was chosen because the
action space is very small, the values can be inspected directly, and the update
equation can be defended clearly.

Several design decisions were refined during development. First, game logic was
moved into \texttt{game.py} so that turn order and round advancement belong to the
environment, not to the learning loop. Second, state representation was studied in
two forms: a full slime-mask representation and a compact reduced representation
using local landing information. Third, reward shaping was adjusted after observing
that excessive fall penalties encourage behavior that is too conservative for a
race. Fourth, a maximum round limit was introduced so that long or stalled games
are counted as timeouts instead of running indefinitely.

The report follows the structure requested by the assignment. Section~II gives the
RL background. Section~III describes the game and the MDP formulation.
Section~IV presents the implementation architecture. Section~V explains the
Q-learning algorithm. Section~VI describes opponents and training. Section~VII
discusses reward design. Section~VIII presents the experimental methodology.
Section~IX analyzes the CSV results. Section~X discusses the main lessons.
Sections~XI--XIII describe limitations, reproducibility, and future work, and
Section~XIV concludes.

\section{Reinforcement Learning Background}

An RL problem is commonly formalized using the assumptions of a Markov Decision
Process (MDP): the agent perceives a finite set of states, has a finite set of
available actions, acts in discrete time, and receives a numerical reward or
penalty from the environment \cite{sanchoQLearning}. At each decision step, the
agent observes a state $s$, selects an action $a$, receives a reward $r$, and
transitions to a new state $s'$. The goal is to learn a policy $\pi$ that maximizes
expected discounted return:

\begin{equation}
G_t = \sum_{k=0}^{\infty} \gamma^k R_{t+k+1},
\end{equation}

where $\gamma \in [0,1]$ is the discount factor.

The Markov property means that the current state contains all information needed
to model the distribution of future states given the chosen action. In practice,
implementations often use a state abstraction. A compact abstraction may violate
the exact Markov property, but it can generalize better when tabular data is
sparse. This tradeoff became central in this project.

Q-learning is a reinforcement learning method in which the agent assigns values to
state-action pairs \cite{sanchoQLearning}. The value $Q(s,a)$ estimates the
long-term usefulness of taking action $a$ in state $s$, combining the immediate
reward with discounted future rewards. The update rule used in this project is the
gradual version of this idea:

\begin{equation}
Q(s,a) \leftarrow Q(s,a) +
\alpha \left[r + \gamma \max_{a'} Q(s',a') - Q(s,a)\right].
\label{eq:qlearning}
\end{equation}

For terminal states, the future value term is omitted. The learning rate $\alpha$
controls how strongly the new sample changes the old estimate, while $\gamma$
controls how much future rewards influence the current value. The main alternative
considered was SARSA, an on-policy method that updates using the action actually
selected next. SARSA could be useful in this game because it may learn more
conservative policies under exploration. However, Q-learning was kept as the main
algorithm because it is simpler, standard, and easier to explain in a first
implementation.

\section{Game and MDP Formulation}

\subsection{Game Rules}

The default game uses a track of length 10. Tiles are indexed from 0 to 9. Both
players start on tile 0, and tile 9 is the goal. Player 0 is the learner and player
1 is controlled by a dummy opponent during training and evaluation. The game ends
immediately when either player reaches the goal.

Every non-terminal, non-start tile may be dry or slimy. At initialization, each
interior tile is independently sampled as slimy with probability 0.5. After both
players have acted in a round, each empty interior tile has probability 0.5 of
flipping to the opposite terrain. Occupied tiles do not flip. This rule matters
because waiting may preserve a currently favorable or unfavorable board, while
moving may expose the agent to future terrain changes.

The available actions are:

\begin{itemize}
\item \texttt{STAY}: remain on the current tile.
\item \texttt{JUMP}: attempt to move one tile forward.
\item \texttt{LONG\_JUMP}: attempt to move two tiles forward.
\end{itemize}

The landing probabilities are shown in Table~\ref{tab:landing}. If a jump fails,
the player returns to tile 0. The goal and start tiles are represented as non-slimy
in the board model.

\begin{table}[h]
\centering
\caption{Landing probabilities}
\label{tab:landing}
\begin{tabular}{lcc}
\toprule
Target terrain & \texttt{JUMP} & \texttt{LONG\_JUMP} \\
\midrule
Dry & 0.90 & 0.70 \\
Slimy & 0.60 & 0.35 \\
\bottomrule
\end{tabular}
\end{table}

\subsection{Action Space}

The action set is:

\begin{equation}
A = \{\texttt{STAY}, \texttt{JUMP}, \texttt{LONG\_JUMP}\}.
\end{equation}

The legal action set depends on the current player position. A one-tile jump is
legal unless the player is already at the goal. A long jump is legal only when the
player is at least two tiles away from the goal. Terminal states have no legal
actions. The implementation enforces legal actions in the game engine rather than
in the learning code, so all players share the same rule validation.

\subsection{Transition Dynamics}

The transition is stochastic for two reasons. First, jump outcomes are sampled by a
Bernoulli trial determined by the target terrain and action type. Second, after
both players act, slime flips are sampled independently for empty interior tiles.
A transition from one learner decision state to the next therefore includes the
learner's landing outcome, the opponent's chosen action and landing outcome, and
the terrain flips at the end of the round.

This structure motivated an important design decision: the learning update is made
after a complete round, not immediately after the learner action. The agent's
reward and next state are computed after the opponent has responded and the terrain
has changed. This better matches the strategic unit of the game. A move is good
not only because it advances the learner, but also because it leaves the learner in
a favorable state after the opponent has had a chance to win or change the race
tempo.

\subsection{State Representation}

Two state encodings were considered.

The first encoding was the full state:

\begin{equation}
s_{\text{full}} = (p_0, p_1, m),
\end{equation}

where $p_0$ is the learner position, $p_1$ is the opponent position, and $m$ is a
binary mask for all slime tiles. This representation is closest to the formal MDP
because it includes the entire terrain configuration. For a track of length $L$,
the approximate upper bound is:

\begin{equation}
|S_{\text{full}}| \approx L^2 2^{L-2}.
\end{equation}

The second encoding was the reduced state:

\begin{equation}
s_{\text{red}} =
(p_0, p_1, t_0^1, t_0^2, t_1^1, t_1^2),
\end{equation}

where $t_i^1$ and $t_i^2$ indicate whether player $i$'s one-tile and two-tile
landing targets are slimy. This state is not a full Markov state because it does
not include terrain farther away. However, it captures the information that most
directly affects the next action and dramatically reduces the number of Q-table
entries. This was not an accidental shortcut: it was an explicit state abstraction
chosen to test whether generalization through abstraction could outperform exact
but sparse tabular learning.

The experiments show that this tradeoff matters. On the 10-tile track, the reduced
state reached 771 known states after 20,000 episodes and performed best. The full
state reached 17,983 known states but performed worse. On a 20-tile track, the
full state reached 5,203,483 known states, showing the curse of dimensionality very
clearly.

\section{Implementation Architecture}

The source code is intentionally small and organized by responsibility. The main
files are listed in Table~\ref{tab:files}. The implementation uses only the Python
standard library for the learning algorithm. No external RL-specific libraries are
used.

\begin{table}[h]
\centering
\caption{Main source files}
\label{tab:files}
\begin{tabular}{ll}
\toprule
File & Responsibility \\
\midrule
\texttt{game.py} & Core rules, actions, turn order, rounds \\
\texttt{dummy\_players.py} & Fixed opponent strategies \\
\texttt{rl.py} & Q-learning, rewards, train/evaluate, JSON policies \\
\texttt{experiments.py} & Reproducible experiment runner \\
\texttt{dummy\_play.py} & Human or dummy terminal game interface \\
\texttt{play\_agent.py} & Human against saved JSON policy \\
\bottomrule
\end{tabular}
\end{table}

\subsection{Game Engine}

The game engine defines the \texttt{Action} enumeration. Each action stores its jump
distance and landing probabilities. Keeping these constants inside the action
definition makes the rule model easy to inspect and avoids duplicating probability
values in the learning code.

The \texttt{Game} object stores the mutable environment state: player positions,
slime positions, the winner, the round number, and the current player. During
development, some round logic was temporarily repeated in \texttt{rl.py}. This was
later refactored so that \texttt{apply\_action()} advances the game automatically.
When player 0 acts, the current player changes to player 1. When player 1 acts,
empty tiles flip, the round counter increases, and the current player resets to
player 0. This design is better because there is only one source of truth for the
rules. The RL code no longer needs to remember when tile flipping should occur.

One consequence of this design is that the game engine also enforces turn order.
Calling \texttt{legal\_actions()} for the wrong player returns an empty tuple, and
calling \texttt{apply\_action()} out of turn raises an error. This prevents invalid
training trajectories from silently entering the Q-table.

\subsection{Policy and Q-Table Storage}

The Q-table is a Python dictionary from state keys to dictionaries of action names
and floating-point values. A typical entry stores values for \texttt{STAY},
\texttt{JUMP}, and \texttt{LONG\_JUMP}. Policies are saved as JSON files containing
the Q-table, the training configuration, and the reward configuration. JSON was
chosen because it is human-readable, easy to generate with the standard library,
and convenient for checking whether policies contain a reasonable number of known
states.

A bug discovered during review was that using \texttt{q\_table or \{\}} in the
learner constructor replaced an intentionally shared empty dictionary with a new
dictionary. This meant the learner appeared to run but did not preserve learning
between episodes. The constructor was corrected to use the provided table whenever
it is not \texttt{None}. This debugging step is important because it shows why
policy size and smoke tests are useful: an empty Q-table can still play by random
tie-breaking, but it is not a learned policy.

\subsection{Human Play Against a Learned Policy}

The project includes a dedicated command-line script, \texttt{play\_agent.py}, for
playing against a saved JSON policy. This script is separate from the training
code because interactive play and experiment execution have different needs. The
experiment runner must simulate many games without input, while the play script
must display the board, ask for human actions, describe moves, and stop cleanly
when the game ends or reaches the maximum round limit.

The script loads the JSON policy file, reads the saved Q-table and training
metadata, and creates a \texttt{QLearningPlayer} attached to a new \texttt{Game}
instance. If the user does not pass a track length or maximum round limit, the
script reuses the values stored in the policy metadata. This avoids an important
class of errors: a policy trained on a 10-tile track should not accidentally be
used on a different track length unless the user explicitly requests it.

By default, the learned policy is assigned to player 1 and the human is assigned
to player 2. This is deliberate. The Q-table was trained from the perspective of
player index 0, so letting the policy move first keeps the interactive behavior
closest to the training setup. The script still provides a \texttt{--human-player}
option, but changing the move order should be interpreted carefully because the
state keys are perspective-based and the policy was not trained symmetrically for
both player indices.

During play, \texttt{play\_agent.py} uses the same board display and input helpers
as the terminal dummy-player game. Human controls are \texttt{s} or \texttt{0} for
stay, \texttt{j} or \texttt{1} for a one-tile jump, and \texttt{l} or \texttt{2}
for a long jump. The learned policy acts greedily with respect to its Q-table: it
does not use exploration during human play. If the current state was not seen
during training, all action values default to zero and the agent breaks ties
randomly among legal actions. This fallback is useful because it keeps the game
playable, but it also makes unseen states easy to recognize during qualitative
testing.

\subsection{Maximum Round Limit}

All training and evaluation functions now accept a maximum round parameter. The
default is 200 for the 10-tile experiments. If a game reaches the limit without a
winner, it is counted as a timeout. During training, a timeout is treated as a
terminal transition for bootstrapping purposes: the final Q-update does not include
future value beyond the cap. This prevents infinite or very long games and makes
experiments reproducible. The CSV output includes wins, losses, timeouts, win
rate, loss rate, timeout rate, and the maximum round value used.

\section{Learning Algorithm}

\subsection{Q-Learning}

The learning algorithm is tabular Q-learning. At the beginning of each episode, a
new game is generated using a deterministic seed based on the base seed and episode
number. The learner controls player 0. The opponent is selected according to the
training-opponent setting. In the main experiments, training uses a mixed opponent
setting that samples from random, cautious, and aggressive strategies.

The learner uses epsilon-greedy exploration. With probability $\epsilon$, it
chooses a random legal action. Otherwise, it chooses a legal action with maximum
Q-value in the current state. This follows the exploration/exploitation tradeoff
described in the Q-learning article: early in training, the agent should try a
wider range of actions because its value estimates are still unreliable; later, it
should rely more on the best values it has learned \cite{sanchoQLearning}. Ties
are broken randomly. Random tie-breaking matters because all unseen actions
initially have value 0. Without random tie-breaking, the agent could develop a
persistent bias toward whichever action appears first in the enumeration.

The training loop can be summarized as follows:

\begin{enumerate}
\item Reset the game and construct the learner and dummy opponent.
\item Encode the learner decision state.
\item Select a learner action using epsilon-greedy exploration.
\item Apply the learner action through the game engine.
\item If the learner has not won, let the opponent choose and apply its action.
\item If the opponent has not won, the game engine flips tiles and advances the
round.
\item Compute the reward for the round.
\item If the game ended or timed out, update without bootstrapping.
\item Otherwise, encode the next state and update using \eqref{eq:qlearning}.
\end{enumerate}

\subsection{Hyperparameters}

The main hyperparameters are shown in Table~\ref{tab:hyper}. These values were
chosen to keep the algorithm simple and stable. The learning rate is moderate:
large enough to respond to new samples, but not so large that one unlucky failed
jump completely overwrites previous experience. The discount factor values future
positioning and eventual wins, which is necessary because most actions do not end
the game immediately.

\begin{table}[h]
\centering
\caption{Main Q-learning hyperparameters}
\label{tab:hyper}
\begin{tabular}{lc}
\toprule
Parameter & Value \\
\midrule
Learning rate $\alpha$ & 0.15 \\
Discount factor $\gamma$ & 0.90 \\
Initial epsilon & 0.30 \\
Final epsilon & 0.02 \\
Epsilon decay & 0.9995 \\
Default track length & 10 \\
Default max rounds & 200 \\
Evaluation games & 1000 per opponent \\
\bottomrule
\end{tabular}
\end{table}

\subsection{Why Q-Learning Was Chosen}

Q-learning was selected because it matches the assignment constraints and the
scale of the game. The state and action spaces are discrete, the action set has
only three elements, and the value table can be saved and inspected. This makes the
method transparent for a defense presentation. It is also simple enough to
implement from scratch correctly.

Other methods were considered. SARSA could make the policy more cautious because
it learns the value of the actually followed exploratory policy. Expected SARSA
could reduce update variance. Double Q-learning could reduce overestimation of
risky long jumps. Value iteration would be possible if all transition
probabilities, opponent policies, and tile flips were enumerated, but that would
shift the project toward planning rather than learning from simulated experience.
For this assignment, Q-learning provides the best balance between originality,
clarity, and sufficient performance.

\section{Dummy Opponents}

The assignment requires simulated opponents because a student cannot manually play
thousands of training games. Four dummy strategies were implemented:

\begin{itemize}
\item \textbf{Random}: chooses uniformly among legal actions.
\item \textbf{Cautious}: \texttt{STAY} if the next tile is slimy; otherwise \texttt{JUMP}.
\item \textbf{Aggressive}: uses \texttt{LONG\_JUMP} whenever legal, otherwise
\texttt{JUMP}.
\item \textbf{Balanced}: uses \texttt{LONG\_JUMP} when the two-tile target is dry and 
\texttt{JUMP} otherwise.
\end{itemize}

The strategies were designed to create different learning pressures. Random is a
weak baseline. Cautious is slow but avoids many obvious risks. Aggressive is fast
but often falls. Balanced is a stronger hand-coded policy because it uses terrain
information while still moving quickly.

The main training setting uses a mixed opponent distribution over random,
cautious, and aggressive opponents. Balanced is primarily used as a challenging
evaluation opponent or as a possible specialized training opponent. This design
tries to avoid overfitting to a single fixed opponent behavior.

\section{Reward Design}

Reward design was one of the most important parts of the project. The game is a
race, so the agent must learn that risk can be useful. A policy that never falls
may still lose because it moves too slowly. A policy that always long-jumps may
also lose because failed jumps reset progress. The reward must therefore encourage
winning and progress without making falling seem worse than losing the whole game.

The shaped reward used in the stored main result summaries is:

\begin{itemize}
\item win: $+10.0$
\item loss: $-10.0$
\item per-round step cost: $-0.02$
\item progress reward: $+0.20$ per tile advanced
\item fall penalty: $-1.0$
\end{itemize}

This reward is a compromise. The terminal values keep winning and losing as the
main objective. The progress term gives intermediate feedback so the learner does
not need to wait until the end of the game to learn that forward movement is
usually useful. The step cost discourages waiting forever. The fall penalty is
large enough to make failed jumps undesirable, but not so large that the agent is
afraid of all risky movement.

During development, harsher reward settings were also considered, including
configurations with a very large win reward and a fall penalty close to half the
loss penalty. Analysis showed that such settings can make the immediate expected
reward of even a dry one-tile jump worse than staying. That is not aligned with the
race objective. This is why the recommended reward interpretation is balanced
shaping, not maximum punishment for falling.

\section{Experimental Methodology}

\subsection{Reproducibility}

Experiments are run through \texttt{src/experiments.py}. The runner accepts the
training budget, training opponent, evaluation opponents, number of evaluation
games, random seed, track length, maximum rounds, and Q-learning hyperparameters.
For every training budget, it saves a policy JSON file and writes CSV/JSON result
summaries under \texttt{outputs/}.

The main command shape is:

\begin{verbatim}
./src/experiments.py --episodes 1000 5000 20000 \
  --eval-games 1000 --seed 42 --max-rounds 200
\end{verbatim}

The 20-tile experiments use a higher maximum round limit:

\begin{verbatim}
./src/experiments.py --episodes 20000 \
  --track-length 20 --max-rounds 5000
\end{verbatim}

\subsection{Compared Conditions}

The experiments compare:

\begin{itemize}
\item reduced-state learning on a 10-tile track;
\item full-state learning on a 10-tile track;
\item reduced-state learning on a 20-tile track;
\item full-state learning on a 20-tile track.
\end{itemize}

For the 10-tile setting, policies were trained for 1,000, 5,000, and 20,000
episodes. For the 20-tile scalability test, policies were trained for 20,000
episodes. Each trained policy was evaluated over 1,000 games against the random,
cautious, and aggressive opponents. A random-player baseline was also evaluated
against the same opponents and seeds.

\subsection{Metrics}

The primary metric is win rate. The CSV files also record loss rate, timeout rate,
known Q-table states, and improvement over the random baseline in percentage
points. The number of known states is not a direct performance metric, but it is
very useful for interpreting sample complexity. If a policy knows millions of
states but performs only slightly above baseline, then the representation is likely
too large for the training budget.

\section{Results and Analysis}

\subsection{Default 10-Tile Track}

Table~\ref{tab:l10reduced} shows the reduced-state results for the default
10-tile game. At 1,000 episodes the agent is not yet reliable. It beats the random
baseline against the cautious opponent by 1.4 percentage points, but it performs
worse against random and aggressive opponents. By 5,000 episodes the policy is
better than baseline in every opponent setting. By 20,000 episodes it achieves a
mean win rate of 0.578, compared with a mean random baseline of 0.411.

\begin{table}[h]
\centering
\caption{Reduced-state results on track length 10}
\label{tab:l10reduced}
\begin{tabular}{lrrrr}
\toprule
Episodes & Opponent & Win & Baseline & Imp. pp \\
\midrule
1,000 & Random & 0.417 & 0.500 & -8.3 \\
1,000 & Cautious & 0.408 & 0.394 & +1.4 \\
1,000 & Aggressive & 0.290 & 0.339 & -4.9 \\
5,000 & Random & 0.558 & 0.500 & +5.8 \\
5,000 & Cautious & 0.572 & 0.394 & +17.8 \\
5,000 & Aggressive & 0.425 & 0.339 & +8.6 \\
20,000 & Random & 0.626 & 0.500 & +12.6 \\
20,000 & Cautious & 0.651 & 0.394 & +25.7 \\
20,000 & Aggressive & 0.456 & 0.339 & +11.7 \\
\bottomrule
\end{tabular}
\end{table}

The improvement is strongest against the cautious opponent. This is expected:
cautious behavior avoids many falls but gives up tempo. A learned policy that is
willing to take calculated jumps can exploit that slowness. The aggressive
opponent remains the hardest of the three in the 10-tile reduced-state setting
because it can win quickly when long jumps succeed.

\subsection{Full State on Track Length 10}

Table~\ref{tab:l10full} shows the full-state results on the same 10-tile track.
The 1,000 episode policy is slightly above baseline on average, but performance
deteriorates at 5,000 and 20,000 episodes. At 20,000 episodes the mean win rate is
0.270, far below the reduced-state mean of 0.578.

\begin{table}[h]
\centering
\caption{Full-state results on track length 10}
\label{tab:l10full}
\begin{tabular}{lrrrr}
\toprule
Episodes & Opponent & Win & Baseline & Imp. pp \\
\midrule
1,000 & Random & 0.521 & 0.500 & +2.1 \\
1,000 & Cautious & 0.436 & 0.394 & +4.2 \\
1,000 & Aggressive & 0.334 & 0.339 & -0.5 \\
5,000 & Random & 0.326 & 0.500 & -17.4 \\
5,000 & Cautious & 0.238 & 0.394 & -15.6 \\
5,000 & Aggressive & 0.140 & 0.339 & -19.9 \\
20,000 & Random & 0.309 & 0.500 & -19.1 \\
20,000 & Cautious & 0.346 & 0.394 & -4.8 \\
20,000 & Aggressive & 0.154 & 0.339 & -18.5 \\
\bottomrule
\end{tabular}
\end{table}

The full-state result is the clearest evidence that exact state information is not
automatically better for tabular learning. The full state has better theoretical
properties, but it separates experiences that the reduced state shares. Two boards
with identical local landing risks but different far-away slime patterns become
different Q-table keys. With only 20,000 episodes, many full-state keys receive too
few updates to estimate values well.

\subsection{Scalability to Track Length 20}

The 20-tile experiments test scalability. Table~\ref{tab:l20} compares the reduced
and full representations after 20,000 episodes with a maximum round limit of 5,000.

\begin{table}[h]
\centering
\caption{Track length 20 results after 20,000 episodes}
\label{tab:l20}
\begin{tabular}{llrrrr}
\toprule
State & Opponent & Win & Baseline & Imp. pp & States \\
\midrule
Reduced & Random & 0.286 & 0.497 & -21.1 & 4,556 \\
Reduced & Cautious & 0.056 & 0.093 & -3.7 & 4,556 \\
Reduced & Aggressive & 0.210 & 0.388 & -17.8 & 4,556 \\
Full & Random & 0.560 & 0.497 & +6.3 & 5,203,483 \\
Full & Cautious & 0.107 & 0.093 & +1.4 & 5,203,483 \\
Full & Aggressive & 0.432 & 0.388 & +4.4 & 5,203,483 \\
\bottomrule
\end{tabular}
\end{table}

The reduced representation performs poorly on the longer track. Local information
is no longer enough: planning over a longer horizon requires more knowledge about
future terrain and accumulated race position. The full representation performs
slightly above baseline, but only after visiting more than five million states.
This is not practical as a final tabular design. The 20-tile results therefore
support two conclusions at the same time: compact local state can work well for the
default game, and full tabular state does not scale cleanly.

\subsection{Known State Counts}

Table~\ref{tab:states} summarizes the state-count tradeoff. The reduced 10-tile
policy saturates at about 771 known states, while the full 20-tile policy reaches
more than five million. The performance improvement from those millions of states
is modest, so the extra table size is not justified for the assignment-sized
solution.

\begin{table}[h]
\centering
\caption{Known Q-table states and mean performance}
\label{tab:states}
\begin{tabular}{lrr}
\toprule
Experiment & Known states & Mean win \\
\midrule
10 tiles, reduced, 20k & 771 & 0.578 \\
10 tiles, full, 20k & 17,983 & 0.270 \\
20 tiles, reduced, 20k & 4,556 & 0.184 \\
20 tiles, full, 20k & 5,203,483 & 0.366 \\
\bottomrule
\end{tabular}
\end{table}

\section{Discussion}

\subsection{What the Agent Learned}

The best default policy learns a useful race strategy. It does not merely avoid
falls; it learns that forward progress is necessary. Its strongest performance is
against cautious opponents, which confirms that it can exploit slow play. Its
weaker performance against aggressive opponents shows that stochastic long-jump
success can still create pressure. This is appropriate for the game: a simple
tabular learner should not be expected to dominate every hand-coded strategy.

\subsection{Reduced State Versus Full State}

The reduced-state result is not a failure of MDP thinking. It is an example of an
engineering tradeoff. A full MDP state is desirable because it contains complete
terrain information. However, a tabular algorithm does not generalize between
different keys. In a small training budget, a compact abstraction can be superior
because it reuses experience across many similar boards. The default 10-tile
results show this very clearly.

For a final assignment implementation, the reduced state is the more appropriate
default. The full state should be presented as an experiment demonstrating the
curse of dimensionality. If the project were extended to longer tracks, a better
solution would be function approximation or a more carefully engineered feature
representation, not a larger raw table.

\subsection{Reward Interpretation}

The reward must be interpreted as part of the model design, not as an objective
truth. The true objective is to win the game. The shaped reward is a tool that
helps temporal-difference learning assign credit before the final outcome. The
fall penalty is especially sensitive. If it is too small, the agent may jump
recklessly. If it is too large, the agent may learn that staying is safer than
moving, even though a race cannot be won by waiting. The selected balanced reward
therefore intentionally values progress and terminal wins more than accident
avoidance.

\subsection{Round Ownership}

Moving round logic into the \texttt{Game} class improved the correctness of the
implementation. Earlier versions of the learning loop manually applied both player
actions and then manually flipped tiles. That duplicated game rules outside the
environment. The final design makes \texttt{apply\_action()} responsible for
advancing turns and rounds. This is cleaner because dummy players, human play, and
RL training all share the same transition mechanics.

\section{Threats to Validity and Limitations}

\subsection{Single Seed}

The CSV files analyzed in this report use seed 42. This makes the experiments
reproducible, but it is not enough to fully characterize stochastic learning. A
stronger final study should repeat each condition with several seeds, for example
7, 42, and 123, and report mean and standard deviation.

\subsection{Policy File Naming}

Some experiment CSV files point to the same policy file names even when the state
representation or track length differs. This means later runs may overwrite earlier
policies. The CSV rows themselves remain useful as recorded results, but final
experiments should include the state mode and track length in policy file names.

\subsection{Interaction Interface}

The terminal interface supports human play against human, dummy, and learned
policy opponents. The learned-policy path is intentionally simple: it loads a JSON
policy, creates the same game environment used during training, and lets the human
enter actions from the terminal. This satisfies the interaction requirement of the
assignment and also provides a qualitative debugging tool. One limitation remains:
because the policy was trained as player 0, the most faithful interactive setup is
for the learned policy to move first. A future version could train mirrored
policies or encode player perspective more explicitly so that either player order
is equally natural.

\subsection{Limited Opponent Models}

The dummy opponents are fixed heuristics. They are useful for controlled
evaluation, but they do not adapt. A learned policy trained only against such
opponents may exploit their patterns without becoming generally strong. Mixed
training reduces this risk but does not remove it.

\subsection{No Function Approximation}

The implementation is intentionally tabular. This makes the project transparent,
but it limits scalability. The 20-tile full-state experiment demonstrates this
clearly: millions of Q-table states are required, and performance improves only
slightly over baseline.

\section{Reproducibility}

The project can be reproduced using the scripts in \texttt{src/}. The most
important command is:

\begin{verbatim}
./src/experiments.py --episodes 1000 5000 20000 \
  --eval-games 1000 --seed 42 --max-rounds 200
\end{verbatim}

Outputs are written to:

\begin{itemize}
\item \texttt{outputs/policies/}: JSON Q-table policies.
\item \texttt{outputs/results/}: CSV and JSON summaries.
\end{itemize}

The CSV files used for this report are:

\begin{itemize}
\item \texttt{experiment\_results\_reduced\_state.csv}
\item \texttt{experiment\_results\_full\_state.csv}
\item \texttt{experiment\_results\_track20\_reduced\_state.csv}
\item \texttt{experiment\_results\_track20\_full\_state.csv}
\end{itemize}

Before final submission, the code and report should be checked so the reward
configuration in \texttt{RewardConfig} matches the configuration saved in the
reported result files. The result summaries used here record the balanced shaped
reward with win $+10$, loss $-10$, step $-0.02$, progress $+0.20$, and fall
$-1.0$.

To play against a saved policy, the command is:

\begin{verbatim}
./src/play_agent.py \
  --policy outputs/policies/q_mixed_20000_seed42.json
\end{verbatim}

The optional \texttt{--human-player}, \texttt{--seed}, \texttt{--track-length}, and
\texttt{--max-rounds} arguments can be used to change the interactive setup. In
the default configuration, the policy is player 1 and the human is player 2.

\section{Future Work}

Several extensions would improve the project. The first is multi-seed evaluation.
The second is adding SARSA or Double Q-learning as comparison algorithms. SARSA
would test whether on-policy learning produces safer behavior; Double Q-learning
would test whether reducing overestimation improves long-jump decisions. The third
is adding state-mode support as a first-class command-line parameter so reduced
and full-state experiments do not require manual code edits. The fourth is adding
more detailed behavioral statistics, such as fall rate, average rounds, action
distribution, long-jump rate on dry targets, and long-jump rate on slimy targets.
Finally, the interactive agent command could be extended with optional explanations
of the Q-values behind each chosen action, which would make it more useful during
the oral defense.

\section{Conclusions}

This project implemented a tabular Q-learning agent for the long-jumps variant of
``The floor is slime!''. The implementation models stochastic landing, changing
terrain, turn order, dummy opponents, training, evaluation, policy saving, and
timeout-aware experiments. The main design lesson is that exact state information
is not always the best practical choice for a tabular method. On the default
10-tile game, the reduced state representation learned a stronger policy than the
full slime-mask representation because it generalized across locally similar
states. After 20,000 episodes, the reduced-state policy achieved a mean win rate
of 0.578 across the evaluated opponents, outperforming the mean random baseline of
0.411.

The experiments also show that scalability is the central limitation. Full-state
learning on a 20-tile track visited more than five million states. This is
technically possible but not elegant or efficient. For larger games, the natural
next step would be function approximation or a more structured feature
representation, which is also the standard direction suggested when state spaces
become too large for explicit Q-tables \cite{sanchoQLearning}. Within the
assignment scope, however, tabular Q-learning with a compact state abstraction is
an appropriate, explainable, and reproducible solution.

\section*{Use of Generative AI}

Generative AI assistance was used during project planning, code review,
implementation refactoring, experiment interpretation, and report drafting. The
assistance was used to compare design alternatives, identify bugs such as Q-table
sharing and state-space growth, and improve the clarity of the final explanation.
All algorithmic decisions, code behavior, and reported results should be reviewed
and understood by the student before submission and defense.

\begin{thebibliography}{4}

\bibitem{sanchoQLearning}
F. Sancho Caparrini, ``Aprendizaje por refuerzo: algoritmo Q Learning,''
\emph{Blog de Fernando Sancho Caparrini}, Universidad de Sevilla. [Online].
Available:
\url{https://www.cs.us.es/~fsancho/Blog/posts/Aprendizaje_por_Refuerzo_Q_Learning}

\bibitem{assignment2025}
A. Riscos Nunez and J. M. Moyano Murillo, ``Artificial Intelligence 2025/26:
Main practical assignment: Reinforcement learning,'' Universidad de Sevilla,
2025.

\bibitem{pythonDocs}
Python Software Foundation, ``Python 3 standard library documentation.'' [Online].
Available: \url{https://docs.python.org/3/library/}

\end{thebibliography}

\end{document}
